\section{Layered codes, joint exchanges and signed constructions}
\label{sec:codes}

These seven results use one geometric construction, with improvements at
different levels. Dimensions 32, 33 and 34 use larger compatible families of
eight-element supports; dimension 39 combines such a family with an explicit
sign code. Dimensions 35, 36 and 37 also change the signs carried by the
supports. The important transition is from deciding whether to retain a whole
bundle to replacing it by a different finite set.

\subsection{Compatibility of the layers}

The Edel--Rains--Sloane construction~\cite{ers} places three coordinate patterns
on the same unit sphere. Let $n\ge m\ge32$, let $T\subset\{0,1\}^m$ be a
binary code, and let $\mathcal B$ be a family of eight-element subsets of
$\{0,\ldots,n-1\}$. A support records the nonzero coordinates of a vector;
the signs specify the vector itself. The dense, middle and root layers are
\begin{align*}
X_T&=\left\{\frac{1}{\sqrt m}
 \bigl((-1)^{c_0},\ldots,(-1)^{c_{m-1}},0,\ldots,0\bigr):c\in T\right\},\\
X_{\mathcal B}&=\left\{\frac{1}{\sqrt8}\sum_{i\in B}\varepsilon_i e_i:
 B\in\mathcal B,\ \varepsilon_i\in\{-1,1\},\ \prod_{i\in B}\varepsilon_i=1\right\},\\
X_R&=\left\{\frac{\pm e_i\pm e_j}{\sqrt2}:0\le i<j<n\right\}.
\end{align*}
Each support carries $2^7=128$ even-parity sign patterns. The layers have
respectively $m$, eight and two nonzero coordinates, so they are disjoint.
Their union contains
\begin{equation}
N=|T|+128|\mathcal B|+2n(n-1)
\label{eq:ers-count}
\end{equation}
points.

\begin{proposition}[Three-layer construction~\cite{ers}]
\label{prop:ers}
If $T$ has minimum Hamming distance at least $\lceil m/4\rceil$ and distinct
supports satisfy $|B\cap B'|\le4$, then the union is a kissing configuration
in dimension $n$.
\end{proposition}
\begin{proof}
Two dense-layer points have product $1-2d_H(c,c')/m\le1/2$.
Distinct even-parity patterns on the same support differ in at least two
positions, giving at most $1-4/8=1/2$. Different supports give at most
$|B\cap B'|/8\le1/2$. Distinct roots either share at most one coordinate
with the same sign, or have the same support and differ in at least one
sign; their product is also at most $1/2$.

There are three remaining cross-layer cases. Dense--middle products are
at most $8/(\sqrt m\sqrt8)=\sqrt{8/m}\le1/2$; dense--root products are
at most $2/(\sqrt m\sqrt2)=\sqrt{2/m}\le1/2$; middle--root products are
at most $2/(\sqrt8\sqrt2)=1/2$. These cover all distinct pairs.
\end{proof}

The geometric problem reduces to finding a large binary code and a large
support family with intersection at most four. The support indicators have
constant weight eight and satisfy
$d_H(1_B,1_{B'})=16-2|B\cap B'|$. Thus the intersection condition is
minimum distance eight in the constant-weight code. Equation~\eqref{eq:ers-count}
converts improvements of such codes into kissing-number lower bounds
\cite{bsss,brouwer,echols}.

\subsection{Joint exchanges of supports}

Inserting a new support generally requires removing old supports with large
intersection. For a single candidate, these deletions may consume the gain.
Several mutually compatible candidates can instead be obstructed by the
same old supports. The search therefore selects additions jointly.

For a candidate $b$, put $F(b)=\{B\in\mathcal B:|B\cap b|>4\}$.
If $Y$ is disjoint from the old family and its internal intersections are
at most four, then
\begin{equation}
\begin{split}
\mathcal B'&=(\mathcal B\setminus F(Y))\cup Y,\qquad
 F(Y)=\bigcup_{b\in Y}F(b),\\
|\mathcal B'|-|\mathcal B|&=|Y|-|F(Y)|.
\end{split}
\label{eq:exchange}
\end{equation}
The cost is a union size, not the sum of the individual conflict counts.
The stored final exchange in dimension 33 removes nine supports and adds
ten; in dimension 34 it removes six and adds seven; an exchange in dimension
37 removes 37 and adds 40. The data contain
the parent families and actual deletions and additions, so these unions
can be reconstructed.

There is also a direct discrete formulation. For a finite candidate set
$\mathcal C$, let $z_b\in\{0,1\}$ select a candidate and let
$u_B\in\{0,1\}$ delete an old support. Impose $z_b+z_{b'}\le1$
for incompatible candidates and $z_b\le u_B$ for $B\in F(b)$, and maximize
\[
\sum_{b\in\mathcal C}z_b-\sum_{B\in\mathcal B}u_B.
\]
An optimal choice deletes only supports actually blocking selected
candidates, so its objective is exactly the gain in
\eqref{eq:exchange}. This identifies the useful search property:
concentrate conflicts on shared old supports, rather than minimizing
each candidate's conflict count separately. Candidate generation and
joint selection are complementary tasks; changing the candidate family
changes the constructions available to this finite optimization problem.

The following parameters give direct three-layer constructions. The
dimension-37 support family was enlarged from 2843 to 2845 and then
modified by a signed replacement, whose final count is given below.
\begin{center}
\begin{minipage}{\linewidth}\centering
\captionof{table}{Parameters of the direct three-layer constructions.}
\label{tab:code-parameters}
\begin{tabular}{rrrrr}
\toprule
$n$ & Dense length $m$ & $|T|$ & $|\mathcal B|$ & Points\\
\midrule
32&32&131072&1676&347584\\
33&32&131072&1803&363968\\
34&32&131072&1960&384196\\
39&39&327680&3383&763668\\
\bottomrule
\end{tabular}
\end{minipage}
\end{center}
Each extra support contributes 128 spherical points. For example,
$131072+128\cdot1676+1984=347584$ in dimension 32.
The gains of 15 and 24 supports in dimensions 33 and 34 give respectively
1920 and 3072 points over the comparison constructions.

The length-32 layer uses the Cheng--Sloane $[32,17,8]$ code
\cite{chengsloane}, whose generator expands to $2^{17}=131072$ words.
Dimension 39 uses $T=\bigcup_{a=1}^{10}(H+r_a)$, where $H$ has dimension
15 and minimum weight 12. Enumeration of the supplied kernel and
representatives gives
\[
\min_{h\in H}\operatorname{wt}(h+r_a+r_b)=10
\]
for all 45 distinct coset pairs. The cosets are disjoint, with
$10\cdot2^{15}=327680$ words and minimum distance ten, as required.
The kernel is obtained by shortening an extended quadratic-residue code;
the programs compare generator row spaces and check the actual distances.

The listed dimension-39 comparison has the same dense-layer size.
Increasing its 3324 supports to 3383 adds $128(3383-3324)=7552$ points.
Recording the support and dense layers separately identifies the finite
objects and their respective contributions.

\subsection{Hadamard transfer and local replacement}

The support count is not the only variable. Each original support carries
all 128 even-parity patterns. Removing a few complete bundles and inserting
a different signed set can fill the affected region more efficiently.
Both the internal products and products against the retained configuration
must be controlled.

\begin{proposition}[Paired Hadamard transfer]
\label{prop:hadamard}
Let $\mathcal A$ be a family of four-element supports with intersections
at most two, transverse to a fixed coordinate pairing. Applying a
Hadamard matrix to each pair in all signed versions of $\mathcal A$
gives $16|\mathcal A|$ weight-eight vectors of squared norm eight and
distinct-pair products at most four.
\end{proposition}
\begin{proof}
Start with four-element supports having pairwise intersections at most
two, and attach all sixteen sign patterns to each. Pair the coordinates
and retain supports meeting each pair in at most one coordinate; call them
transverse. Apply
\[
H=\begin{pmatrix}1&1\\1&-1\end{pmatrix}
\]
to each coordinate pair. Both $(\varepsilon,0)$ and $(0,\varepsilon)$
become two nonzero coordinates, so weight four becomes weight eight.
Distinct input vectors have product at most two: different supports
satisfy the intersection condition, and distinct sign patterns on the same
support differ in at least one position. Since $H^{\mathsf T}H=2I$,
the images have squared norm eight and pairwise products at most four.
Invertibility preserves their number.
\end{proof}

After division by $\sqrt8$, these vectors form an internally compatible
spherical set $Q$. Intersections at most four with each retained support
ensure compatibility with the retained middle layer. The dense and root
cross bounds remain $\sqrt{8/32}=1/2$ and $1/2$, independently of the
parity of the new signs. This is why replacements need not consist of
complete even-parity bundles.

The free-zone replacement principle appears in
Takhanov--Yun~\cite[Section V, Definition 5 and Lemma 4]{signed-johnson}
for signed weight-four codes. Here it is applied to weight-eight
regions inside the ERS middle layer. The choices of region, transverse
supports and paired transform give the replacements below.
In dimensions 35 and 36, boundary compatibility can be secured before
choosing $Q$. Take a coordinate region $W$, remove the old supports
contained in it, and require $|B\cap W|\le4$ for every retained support.
Any new weight-eight vector on $W$ is then compatible with the retained
middle layer. Only its local internal products remain to be controlled.
The regions are $W=\{12,\ldots,22\}$ and $W=\{12,\ldots,23\}$,
respectively, with zero-based coordinates. They contain one and three
old supports. The deletion costs are fixed at 128 and 384 points,
turning a problem involving the whole configuration into a smaller
independent construction problem.

Direct search among signed weight-eight vectors in eleven and twelve
coordinates has respectively $2^8\binom{11}{8}=42240$ and
$2^8\binom{12}{8}=126720$ candidates. Weight-four transfer replaces this
search with the choice of a combinatorial source and a coordinate pairing,
while preserving an algebraic proof of the inner-product bounds.
Dimension 35 starts with a Steiner quadruple system $S(3,4,10)$:
each triple belongs to exactly one block, giving $\binom{10}{3}/4=30$
blocks with pairwise intersections at most two. Enumerating the 945
perfect matchings on ten coordinates yields sixteen transverse blocks.
For dimension 36 we use the 51-block construction of
Kalbfleisch--Stanton~\cite[Theorem 7]{kalbfleisch-stanton}, also described
in~\cite[Section V, Construction 1]{signed-johnson}.
Split twelve coordinates into two sets of six and
use the same five one-factors of $K_6$ in each half. For every colour,
the nine choices of a left and a right edge give a four-element block,
producing 45 crossing blocks. Add, in each half, the complements of
the three edges of the first colour to obtain six further blocks.

The resulting 51 blocks still have intersections at most two.
Same-colour crossing blocks sharing one edge intersect in two points;
different colours share at most one endpoint in each half. Internal
blocks in the same half intersect in two points; those in different
halves are disjoint. An internal block meets a crossing block in at most
two points.

\begin{proposition}[A 576-point local code]
\label{prop:local36}
The paired Hadamard image of 36 transverse blocks from the two
one-factorizations is a 576-point weight-eight code in twelve coordinates.
\end{proposition}
\begin{proof}
Choose the three edges of one colour as the coordinate pairs in each
half. Each edge of the other four colours meets every chosen pair at most
once. Its union with a same-colour edge from the other half is therefore
transverse. There are $4\cdot3\cdot3=36$ such blocks; their intersections
are at most two by the preceding argument. Proposition~\ref{prop:hadamard}
gives $36\cdot16=576$ points. The six internal blocks are unnecessary
for this transverse subfamily.
\end{proof}

In the supplied local coordinate order, the chosen pairs are
\[
(0,1),\ (2,4),\ (3,5),\ (6,7),\ (8,10),\ (9,11).
\]
They form the same colour in
the two halves. The direct construction with these pairs regenerates
exactly the stored 576 vectors. The accompanying program also enumerates
all 945 and 10395 matchings of the two source codes, recovering the
256-point and 576-point sets. Thus the existence proof is explicit,
while the exhaustive calculation records the finite choices available
in this model.

The three constructions select 16, 36 and 32 transverse supports,
producing 256, 576 and 512 new points.
The compatible region in dimension 35 has eleven coordinates; ten form
five pairs and the remaining coordinate is zero. Dimensions 36 and 37
use twelve coordinates each. The explicit pairs and transverse supports
regenerate all 256, 576 and 512 signed points for comparison with the final sets.
\begin{center}
\begin{minipage}{\linewidth}\centering
\captionof{table}{Signed replacements. The support count precedes the deletion of bundles.}
\label{tab:signed-replacements}
\begin{tabular}{rrrrr}
\toprule
Dimension & Original supports & Bundles removed & Points added & Net gain\\
\midrule
35&2157&1&256&128\\
36&2742&3&576&192\\
37&2845&3&512&128\\
\bottomrule
\end{tabular}
\end{minipage}
\end{center}
Combining the support choices and replacements gives
\begin{align*}
K(35)&\ge131072+128(2157-1)+256+2(35)(34)=409676,\\
K(36)&\ge131072+128(2742-3)+576+2(36)(35)=484760,\\
K(37)&\ge131072+128(2845-3)+512+2(37)(36)=498024.
\end{align*}
Relative to the comparison with 2832 supports, dimension 37 gains 1664
points from thirteen extra supports and 128 from the replacement, totalling
1792. Dimensions 35 and 36 keep their original support counts; their
gains come entirely from the signed choices.

The local choice in dimension 35 also attains an exact capacity bound.
Let $M(11)$ be the largest number of weight-eight vectors in
$\{0,\pm1\}^{11}$ with distinct-pair products at most four.
Fix a full sign string $\sigma\in\{\pm1\}^{11}$. Two vectors agreeing
with $\sigma$ on their supports would intersect in at least five
positions and have product at least five. Thus each sign string
extends at most one selected vector. Each vector, in turn, has eight
extensions on its three zero coordinates. Counting gives
\[
8|Q|\le2^{11},\qquad M(11)\le256.
\]
The 256-point construction attains equality, so $M(11)=256$.
This explains the complete use of the eleven-coordinate local region.
In twelve coordinates, two eight-element supports can intersect in
only four positions, allowing a sign string to extend several vectors;
this is the geometric opening used by the larger local code in
dimension 36.

The finite checks use the coordinate order, pairs and four-subsets
specified in the data. In dimension 37, a support string numbered from
one at the left has coordinate $i$ at integer-mask bit $37-i$.
The program regenerates all 512 points from 32 transverse supports and
six coordinate pairs, checks equality with the stored replacement, and
checks intersections with every retained support. The origin, count and
geometric conditions are thereby connected through the same input.

These seven results treat the support family and its sign patterns as
coupled choices. The next construction changes another part of the
geometry: equatorial directions left available by a large lattice lifting.
